\chapter{The statements}\label{chap:statements}

\section*{At a glance}
\begin{itemize}
\item \lead{Goal} State the eleven results that make up the verification, as definitions of type
  \code{Prop} that \code{proof/spec.pin} pins.
\item \lead{Main results} The three milestones \code{erases\_correct}, \code{erasePure\_erases} and
  \code{srcEval\_sound}, and the lemmas \code{collectDeps\_ok\_of}, \code{collectDeps\_sub}, \code{markOf\_det}, \code{markOf\_exists},
  \code{sound\_nonvacuous} and \code{clCert\_proof\_pt}; the final theorem \code{erase\_correct} and its
  first-order corollary are in Chapter~\ref{chap:final}.
\item \lead{Status} \stProved{} the eleven definitions, sorry-free and pinned (checks C12, C12b), and
  \hl{all eleven theorems they state}; the checks compare each theorem's type with its definition.
\item \lead{Nodes} \bpnodes{chap:statements}
\item \lead{Lean files} \leanfile{proof/EraseProof/Statements.lean}.
\end{itemize}

Each statement is a definition node that cites the \code{\_statement} definition. Beside it, a
node states the theorem the definition is the type of.

\begin{tabular}{p{0.28\linewidth}p{0.64\linewidth}}
\toprule
Condition & Meaning \\
\midrule
\code{NoMarkKeys}, \code{NoMarkKeysDecls} & no \code{mdata} with the marks' key in the input, resp. in the types and values of the closure (\code{hclean}; acceptance gives the second) \\
\code{ValuesMarkedOf} & every value of a constant of the closure that the traversal erases (\code{ErasedValue}) has a marking in the empty context (\code{erasePure\_erases}; \code{ProgCert} gives it to the final theorem) \\
\code{SubEnv}, \code{NoPreDefs} & the closure is part of the accepted program; no closure declaration has a compiled pre-definition \\
\code{ClosedDecls} & no free variable in the closure's types and values (\code{markOf\_det}, \code{markOf\_exists}) \\
\code{CtorsListed}, \code{NatShaped}, \code{NatCtors}, \code{RecsShaped} & shape of constructors, \code{Nat} and recursors (\code{erases\_correct}; Definition~\ref{def:EraseProof-CtorsListed}) \\
\code{RulesCanonical}, \code{TablesCanonical} & stored recursor rules have the generated shape and select by constructor; projection tables name their constructor (parts of \code{ConLecheAccepts}) \\
\bottomrule
\end{tabular}

\section{Milestones}\label{sec:statements-milestones}

\begin{definition}[Simulation: statement]
  \label{def:EraseProof-erases_correct-statement}
  \lean{EraseProof.erases_correct_statement}
  \leanok
  \srcloc{proof/EraseProof/Statements.lean}{28}
  \lead{In short} The statement of \code{erases\_correct}, a definition of type \code{Prop}.
\end{definition}

\begin{theorem}[Simulation]
  \label{thm:EraseProof-erases_correct}
  \lean{EraseProof.erases_correct}
  \leanok
  \uses{def:EraseProof-erases_correct-statement}
  \srcloc{proof/EraseProof/Simulation/Main.lean}{23}
  \lead{In short} \hl{A marked source term that evaluates erases to a $\lambda_\Box$ term that evaluates to an erasure of the value}, with no typing premise and no certificate.
  \begin{itemize}
  \item \lead{Given} an erasure of the marked term with erased dependencies; \code{KernameInj}, \code{BlockNamesInj}, \code{CtorsListed}, \code{NatShaped}, \code{NatCtors}, \code{RecsShaped} on the closure; a source evaluation.
  \item \lead{Then} the value has an erasure that the $\lambda_\Box$ term evaluates to.
  \end{itemize}
  \lead{Reference} \code{erases\_correct} (\code{MR E/ErasureCorrectness.v:51}); DV-32.
\end{theorem}
\begin{proof}
  \uses{def:EraseProof-erases_correct-statement, lem:EraseProof-SimCases-erases-correct-beta, lem:EraseProof-SimCases-erases-correct-ctor, lem:EraseProof-SimCases-erases-correct-quotMk, lem:EraseProof-SimCases-erases-correct-recIota}
  \leanok
  \lead{Idea} Induction on the source evaluation; each of its rules is closed by one case lemma (Chapter~\ref{chap:simulation}).
\end{proof}

\begin{definition}[Pure traversal: statement]
  \label{def:EraseProof-erasePure_erases-statement}
  \lean{EraseProof.erasePure_erases_statement, EraseProof.ValuesMarkedOf}
  \leanok
  \srcloc{proof/EraseProof/Statements.lean}{61}
  \lead{In short} The statement of \code{erasePure\_erases}, a definition of type \code{Prop}; \code{ValuesMarkedOf}: every value of a constant of the closure that the traversal erases has a marking in the empty context.
\end{definition}

\begin{theorem}[Pure traversal]
  \label{thm:EraseProof-erasePure_erases}
  \lean{EraseProof.erasePure_erases}
  \leanok
  \uses{def:EraseProof-erasePure_erases-statement}
  \srcloc{proof/EraseProof/Core/Assemble.lean}{159}
  \lead{In short} \hl{The output of the pure traversal on the closure the fragment test returns is an erasure of the marking of the input}, with erased dependencies, in a closed environment.
  \begin{itemize}
  \item \lead{Given} an accepted program of which the closure is part; no compiled pre-definition in the closure; a marking of every value of the closure that the traversal erases and of the input.
  \item \lead{Then} the program is untyped, its main term erases the marking, its environment erases the dependencies and is closed.
  \end{itemize}
  \lead{Reference} \code{erases\_erase} (\code{MR E/ErasureFunction.v:1228}); DV-32, DV-35.
\end{theorem}
\begin{proof}
  \uses{def:EraseProof-erasePure_erases-statement, lem:EraseProof-ConLecheAccepts-namesClosed, lem:EraseProof-DepClosed, lem:EraseProof-Steps-ExprSpec, lem:EraseProof-Steps-erasePure_erases_of, lem:EraseProof-coreEnv_of, lem:EraseProof-ctorPropBoxed_of, lem:EraseProof-collectDeps-spec}
  \leanok
  \lead{Idea} The environment facts and the oracle's facts follow from acceptance; the traversal's invariant at every fuel (Chapter~\ref{chap:core}) gives the conclusion at the empty context.
\end{proof}

\begin{definition}[Meaning of evaluation: statement]
  \label{def:EraseProof-srcEval_sound-statement}
  \lean{EraseProof.srcEval_sound_statement}
  \leanok
  \srcloc{proof/EraseProof/Statements.lean}{81}
  \lead{In short} The statement of \code{srcEval\_sound}, a definition of type \code{Prop}.
\end{definition}

\begin{theorem}[Meaning of evaluation]
  \label{thm:EraseProof-srcEval_sound}
  \lean{EraseProof.srcEval_sound}
  \leanok
  \uses{def:EraseProof-srcEval_sound-statement}
  \srcloc{proof/EraseProof/Semantics/Sound.lean}{762}
  \lead{In short} The marked source evaluation preserves meaning in every con-leche model of the accepted program, except where the input denotes the proof point.
  \begin{itemize}
  \item \lead{Given} con-leche accepts the program and types the input; the input and the closure carry no mark key; the input has a marking that evaluates.
  \end{itemize}
\end{theorem}
\begin{proof}
  \uses{def:EraseProof-srcEval_sound-statement, lem:EraseProof-Sound-of-semSim, lem:EraseProof-semSim_of_srcEval, lem:EraseProof-ConLecheAccepts-namesClosed, lem:EraseProof-Resolve-of_read, lem:EraseProof-markOf-read}
  \leanok
  \lead{Idea} Induction on the marked evaluation with the invariant that the marked term reads in con-leche; each rule follows from the step lemma of Chapter~\ref{chap:semantics}.
\end{proof}

\begin{definition}[Lower bound of the fragment test: statement]
  \label{def:EraseProof-collectDeps_ok_of-statement}
  \lean{EraseProof.collectDeps_ok_of_statement}
  \leanok
  \srcloc{proof/EraseProof/Statements.lean}{44}
  \lead{In short} The statement of \code{collectDeps\_ok\_of}, a definition of type \code{Prop}.
\end{definition}

\begin{lemma}[Lower bound of the fragment test]
  \label{lem:EraseProof-collectDeps_ok_of}
  \lean{EraseProof.collectDeps_ok_of}
  \leanok
  \uses{def:EraseProof-collectDeps_ok_of-statement}
  \srcloc{proof/EraseProof/Core/CollectOk.lean}{1056}
  \lead{In short} A program that con-leche accepts, with an input con-leche types and a closure inside the fragment, is accepted by the fragment test.
  \begin{itemize}
  \item \lead{Why} it shows that \code{hfrag} does not make the final theorem empty.
  \end{itemize}
\end{lemma}
\begin{proof}
  \uses{def:EraseProof-collectDeps_ok_of-statement, lem:EraseProof-ConLecheAccepts-stored, lem:EraseProof-collectDeps-spec, lem:EraseProof-CollectOk-Frag-closure_ok_of}
  \leanok
  \lead{Idea} The fragment test succeeds when every declaration of the closure is in the fragment and the fuel suffices.
\end{proof}

\section{Lemmas}\label{sec:statements-lemmas}

\begin{definition}[Meaning has content: statement]
  \label{def:EraseProof-sound_nonvacuous-statement}
  \lean{EraseProof.sound_nonvacuous_statement}
  \leanok
  \srcloc{proof/EraseProof/Statements.lean}{211}
  \lead{In short} The statement of \code{sound\_nonvacuous}, a definition of type \code{Prop}.
\end{definition}

\begin{lemma}[Meaning has content]
  \label{lem:EraseProof-sound_nonvacuous}
  \lean{EraseProof.sound_nonvacuous, EraseProof.looseBVarsBounded_mkAppN_head}
  \leanok
  \uses{def:EraseProof-sound_nonvacuous-statement}
  \srcloc{proof/EraseProof/Semantics/NonVacuous.lean}{31}
  \lead{In short} If con-leche types an application of the input with a sort, the input has a reading that is not the proof point in every model.
\end{lemma}
\begin{proof}
  \uses{def:EraseProof-sound_nonvacuous-statement}
  \leanok
  \lead{Idea} A member of a sort is a type, never the proof point, and applying the proof point gives the proof point.
\end{proof}

\begin{definition}[Proof certificate: statement]
  \label{def:EraseProof-clCert_proof_pt-statement}
  \lean{EraseProof.clCert_proof_pt_statement}
  \leanok
  \srcloc{proof/EraseProof/Statements.lean}{228}
  \lead{In short} The statement of \code{clCert\_proof\_pt}, a definition of type \code{Prop}.
\end{definition}

\begin{lemma}[Proof certificate]
  \label{lem:EraseProof-clCert_proof_pt}
  \lean{EraseProof.clCert_proof_pt, EraseProof.proof_ptAt, EraseProof.openLams_pt}
  \leanok
  \uses{def:EraseProof-clCert_proof_pt-statement}
  \srcloc{proof/EraseProof/Semantics/Cert.lean}{103}
  \lead{In short} A subterm certified as a proof denotes the proof point under every valuation of its binders.
\end{lemma}
\begin{proof}
  \uses{def:EraseProof-clCert_proof_pt-statement}
  \leanok
  \lead{Idea} The claims for inference on the subterm and its type, and the fact that a member of a proposition is the proof point.
\end{proof}

\begin{definition}[Marking is a function: statement]
  \label{def:EraseProof-markOf_det-statement}
  \lean{EraseProof.markOf_det_statement}
  \leanok
  \srcloc{proof/EraseProof/Statements.lean}{251}
  \lead{In short} The statement of \code{markOf\_det}, a definition of type \code{Prop}.
\end{definition}

\begin{lemma}[Marking is a function]
  \label{lem:EraseProof-markOf_det}
  \lean{EraseProof.markOf_det}
  \leanok
  \uses{def:EraseProof-markOf_det-statement}
  \srcloc{proof/EraseProof/Oracle/Det.lean}{626}
  \lead{In short} For a closure without free variables (\code{ClosedDecls}) and a closed input, the marking at no locals is a function of the term and the configuration.
\end{lemma}
\begin{proof}
  \uses{def:EraseProof-markOf_det-statement, lem:EraseProof-MarkOf-det_ren, lem:EraseProof-Oracle-Rename}
  \leanok
  \lead{Idea} The marking commutes with renaming (Section~\ref{sec:core-oracle}); at no locals the renaming is the identity.
\end{proof}

\begin{definition}[Marking exists: statement]
  \label{def:EraseProof-markOf_exists-statement}
  \lean{EraseProof.markOf_exists_statement}
  \leanok
  \srcloc{proof/EraseProof/Statements.lean}{261}
  \lead{In short} The statement of \code{markOf\_exists}, a definition of type \code{Prop}: \code{hmark} follows from \code{hfrag} and \code{hrun}.
\end{definition}

\begin{lemma}[Marking exists]
  \label{lem:EraseProof-markOf_exists}
  \lean{EraseProof.markOf_exists}
  \leanok
  \uses{def:EraseProof-markOf_exists-statement, lem:EraseProof-Steps-existsSpec_of, lem:EraseProof-coreEnv_of}
  \srcloc{proof/EraseProof/Core/Exists.lean}{138}
  \lead{In short} \hl{Wherever the pure traversal erases an input without mark keys}, over a closed closure that is part of an accepted program, the input has a marking: \code{hmark} excludes no program.
\end{lemma}
\begin{proof}
  \uses{def:EraseProof-markOf_exists-statement, lem:EraseProof-Steps-existsSpec_of, lem:EraseProof-coreEnv_of}
  \leanok
  \lead{Idea} Induction on the fuel: each successful step of \code{visitExpr} yields the marking rule of its case (Section~\ref{sec:core-exists}).
\end{proof}

\begin{definition}[Closure inside the program: statement]
  \label{def:EraseProof-collectDeps_sub-statement}
  \lean{EraseProof.collectDeps_sub_statement}
  \leanok
  \srcloc{proof/EraseProof/Statements.lean}{274}
  \lead{In short} The statement of \code{collectDeps\_sub}, a definition of type \code{Prop}.
\end{definition}

\begin{lemma}[Closure inside the program]
  \label{lem:EraseProof-collectDeps_sub}
  \lean{EraseProof.collectDeps_sub}
  \leanok
  \uses{def:EraseProof-collectDeps_sub-statement}
  \srcloc{proof/EraseProof/Core/Collect.lean}{512}
  \lead{In short} The closure the fragment test returns is part of the program con-leche accepted.
\end{lemma}
\begin{proof}
  \uses{def:EraseProof-collectDeps_sub-statement, lem:EraseProof-CollectSub-closure-sub}
  \leanok
  \lead{Idea} Every constant the fragment test reads comes from a record that the acceptance check covers (Chapter~\ref{chap:core}).
\end{proof}

\section*{Caveats}
\begin{itemize}
\item A statement definition is a type; the checks C12 and C12b compare each theorem's type with it and pin what it reads.
\item The statements read the shipping functions \code{collectDeps}, \code{erasePure} and
  \code{eraseEntry} by their logical definitions.
\item The hypotheses of \code{erase\_correct} are discharged on concrete programs by the tests of
  the proof package (Chapter~\ref{chap:tests}); for the programs with an inductive type, the fold equation of \code{hacc} and a few certificates rest on compiled evaluation
  (Section~\ref{sec:trust-gap}).
\end{itemize}
